%=============================================================================
% OPTIONAL SUBSECTION --- the corrected mechanistic arm
%
% NOT wired into main.tex. To include it, add after the results section:
%     %=============================================================================
% OPTIONAL SUBSECTION --- the corrected mechanistic arm
%
% NOT wired into main.tex. To include it, add after the results section:
%     %=============================================================================
% OPTIONAL SUBSECTION --- the corrected mechanistic arm
%
% NOT wired into main.tex. To include it, add after the results section:
%     %=============================================================================
% OPTIONAL SUBSECTION --- the corrected mechanistic arm
%
% NOT wired into main.tex. To include it, add after the results section:
%     \input{sections/05b_mechanism_corrected}
%
% Protocol note: these numbers come from the *corrected* case series (559 weeks,
% `rebuilt`) under nine disjoint origins, not the 459-week array and three
% origins used in the rest of the paper. The first paragraph says so explicitly,
% because a reviewer will otherwise ask. Every figure is reproducible from
% analysis/_build/ and is logged as EXP-032 to EXP-037.
%=============================================================================

\subsection{The mechanism, specified correctly}
\label{sec:mechanism-corrected}

A natural objection to a negative result about mechanistic structure is that the
mechanism was built wrong. We tested that directly. The arm reported here runs on
the date-ordered case series (559 weeks, 25 districts) under nine disjoint
rolling origins spanning 0.50--0.90 of the record, three seeds each; the protocol
is otherwise that of \S\ref{sec:setup}. We state the change of dataset and fold
count plainly because the comparison inside this subsection is internally paired
and should not be read across to Figure~\ref{fig:results}.

Four defects in the compartmental arm were found and corrected, each tested in
isolation: compartments initialised from weekly flows rather than standing stocks,
which caps the attainable prediction at roughly half the observed count;
subtraction of ten years of cumulative cases from the susceptible pool, which
assumes lifelong immunity to a four-serotype disease; a force of infection with
no dependence on current prevalence, so the model contains no epidemic growth;
and a reporting rate fixed at $1/11$ against a published range of $1/2.5$ to
$1/30$. Corrected, the arm improves from test RMSE 50.13 to 29.86, better on
\textbf{9 of 9} origins, $p = 0.0039$ by a sign-flip permutation test clustered
by origin. The improvement replicates on a second construction of the dataset.

\textbf{The compartmental layer earns its place.} The decisive control is the
identical encoder, loss, normalisation and training budget with the simulator
removed from the prediction path, the network predicting counts directly:
\textbf{30.04 with the mechanism against 49.45 without}, better on 8 of 9
origins, $p = 0.0078$ --- a 39\% reduction in error attributable to the structure
alone. It holds at every horizon tested, by 18.2 RMSE at three weeks, 12.1 at six
and 9.0 at twelve, with $p \le 0.012$ in each case. The no-mechanism arm scores
38.23 on validation and 49.45 on test: it fits the periods it has seen and
extrapolates badly, and the compartmental layer removes exactly that freedom.

\textbf{The graph does not.} Freezing the encoder and replacing it with a single
learned transmission rate changes nothing ($p = 0.97$). Replacing border
adjacency with a population-gravity graph does not help either, and the learned
import weight $\alpha$ converges to exactly zero in all 108 runs, for both
graphs: given a free non-negative parameter for its neighbours' prevalence, the
model declines it.

\textbf{And the ceiling holds.} The corrected arm remains level with persistence
rather than ahead of it ($+1.3$ RMSE over nine origins, $p = 0.92$; the best
variant reaches 28.74 against the floor's 28.54). The mechanism for that is
measurable: the corrected forecast correlates with last-value persistence at
$r = 0.97$--$0.99$, differing by 12--15\% of the level, and a learned convex
blend of the two settles at $w = 0.53$--$0.66$ on persistence while changing the
score by 0.04. Trained on squared error, the mechanism rediscovers persistence
and adds a correction worth one to two RMSE, which is the whole margin available.

A fitted $R_{\mathrm{eff}} = \beta S / \gamma$ of $0.40$--$0.44$ explains why.
At a three-week horizon the forecast is dominated by the incubation and reporting
pipeline draining --- people already infected at the origin --- rather than by new
transmission. Two consequences follow, and both are observed. Forecasting further
ahead does not help: persistence stays ahead at six and twelve weeks and the two
degrade at indistinguishable rates, $1.56\times$ against $1.58\times$ from three
weeks to twelve. And an intervention covariate cannot repair the one window
where the record collapses for a non-epidemiological reason. In 2020 national
notifications fall from 3057 to 299 a week under movement restrictions. We gave
the model the OxCGRT stringency index three ways, once with Google workplace
mobility alongside, over 189 further matched evaluations on the same nine origins
and seeds. Learned from the data it cannot act where it is needed: in all 24 runs
whose training data hold no lockdown week both fitted weights stay at exactly
zero, and pooled over nine origins the learned arms are marginally worse than the
cases-only model, 30.05 against 30.06 and 30.07 at Holm $p = 1.0$, by margins
under the pipeline's own run-to-run spread of 0.109 RMSE. Imposing it as a prior
removes that objection: $\beta \mapsto \beta(1 - \sigma/100)^{\gamma}$ with
$\gamma$ fixed in advance at $0.5$, $1$ and $2$ needs no training exposure, since
restricting movement reducing contact is knowledge from outside this series. It
behaves as the mechanism predicts and still loses. On the 2020 collapse it helps,
monotonically in $\gamma$ --- 63.45 to 63.28, 63.12 and 62.80, over-prediction
bias falling from $+21.4$ to $+19.9$; on the recovery window straight after it is
destroyed --- 11.54 to 13.45, 15.61 and 18.34, bias flipping from $+0.1$ to
$-6.4$. Pooled it is significantly worse, 30.59, 31.28 and 32.43 against 30.07,
better on 1 of 9 origins at $\gamma \ge 1$ ($p = 0.047$, $p = 0.031$).
Stringency stayed high into 2022 while notifications were already rebounding, so a
multiplier that correctly suppresses transmission in April 2020 goes on
suppressing it for two years too long. A correctly signed, correctly placed,
externally supplied intervention term buys 0.65 RMSE on the window it was built
for and costs 6.81 on the next.

\textbf{What the mechanism is for.} Two roles survive. It constrains: the 39\%
reduction above is obtained by forbidding the network to emit counts that no
compartmental flow could produce. And it supplies quantities a point forecast
cannot --- a transmission rate, an implied $R_{\mathrm{eff}}$, a calibrated
reporting rate and compartment states, one of which is the largest single
contributor to ranking which districts breach an alert threshold
(\S\ref{sec:results}). Neither role asks the mechanism to predict the count, and
that is the distinction this paper draws.

%
% Protocol note: these numbers come from the *corrected* case series (559 weeks,
% `rebuilt`) under nine disjoint origins, not the 459-week array and three
% origins used in the rest of the paper. The first paragraph says so explicitly,
% because a reviewer will otherwise ask. Every figure is reproducible from
% analysis/_build/ and is logged as EXP-032 to EXP-037.
%=============================================================================

\subsection{The mechanism, specified correctly}
\label{sec:mechanism-corrected}

A natural objection to a negative result about mechanistic structure is that the
mechanism was built wrong. We tested that directly. The arm reported here runs on
the date-ordered case series (559 weeks, 25 districts) under nine disjoint
rolling origins spanning 0.50--0.90 of the record, three seeds each; the protocol
is otherwise that of \S\ref{sec:setup}. We state the change of dataset and fold
count plainly because the comparison inside this subsection is internally paired
and should not be read across to Figure~\ref{fig:results}.

Four defects in the compartmental arm were found and corrected, each tested in
isolation: compartments initialised from weekly flows rather than standing stocks,
which caps the attainable prediction at roughly half the observed count;
subtraction of ten years of cumulative cases from the susceptible pool, which
assumes lifelong immunity to a four-serotype disease; a force of infection with
no dependence on current prevalence, so the model contains no epidemic growth;
and a reporting rate fixed at $1/11$ against a published range of $1/2.5$ to
$1/30$. Corrected, the arm improves from test RMSE 50.13 to 29.86, better on
\textbf{9 of 9} origins, $p = 0.0039$ by a sign-flip permutation test clustered
by origin. The improvement replicates on a second construction of the dataset.

\textbf{The compartmental layer earns its place.} The decisive control is the
identical encoder, loss, normalisation and training budget with the simulator
removed from the prediction path, the network predicting counts directly:
\textbf{30.04 with the mechanism against 49.45 without}, better on 8 of 9
origins, $p = 0.0078$ --- a 39\% reduction in error attributable to the structure
alone. It holds at every horizon tested, by 18.2 RMSE at three weeks, 12.1 at six
and 9.0 at twelve, with $p \le 0.012$ in each case. The no-mechanism arm scores
38.23 on validation and 49.45 on test: it fits the periods it has seen and
extrapolates badly, and the compartmental layer removes exactly that freedom.

\textbf{The graph does not.} Freezing the encoder and replacing it with a single
learned transmission rate changes nothing ($p = 0.97$). Replacing border
adjacency with a population-gravity graph does not help either, and the learned
import weight $\alpha$ converges to exactly zero in all 108 runs, for both
graphs: given a free non-negative parameter for its neighbours' prevalence, the
model declines it.

\textbf{And the ceiling holds.} The corrected arm remains level with persistence
rather than ahead of it ($+1.3$ RMSE over nine origins, $p = 0.92$; the best
variant reaches 28.74 against the floor's 28.54). The mechanism for that is
measurable: the corrected forecast correlates with last-value persistence at
$r = 0.97$--$0.99$, differing by 12--15\% of the level, and a learned convex
blend of the two settles at $w = 0.53$--$0.66$ on persistence while changing the
score by 0.04. Trained on squared error, the mechanism rediscovers persistence
and adds a correction worth one to two RMSE, which is the whole margin available.

A fitted $R_{\mathrm{eff}} = \beta S / \gamma$ of $0.40$--$0.44$ explains why.
At a three-week horizon the forecast is dominated by the incubation and reporting
pipeline draining --- people already infected at the origin --- rather than by new
transmission. Two consequences follow, and both are observed. Forecasting further
ahead does not help: persistence stays ahead at six and twelve weeks and the two
degrade at indistinguishable rates, $1.56\times$ against $1.58\times$ from three
weeks to twelve. And an intervention covariate cannot repair the one window
where the record collapses for a non-epidemiological reason. In 2020 national
notifications fall from 3057 to 299 a week under movement restrictions. We gave
the model the OxCGRT stringency index three ways, once with Google workplace
mobility alongside, over 189 further matched evaluations on the same nine origins
and seeds. Learned from the data it cannot act where it is needed: in all 24 runs
whose training data hold no lockdown week both fitted weights stay at exactly
zero, and pooled over nine origins the learned arms are marginally worse than the
cases-only model, 30.05 against 30.06 and 30.07 at Holm $p = 1.0$, by margins
under the pipeline's own run-to-run spread of 0.109 RMSE. Imposing it as a prior
removes that objection: $\beta \mapsto \beta(1 - \sigma/100)^{\gamma}$ with
$\gamma$ fixed in advance at $0.5$, $1$ and $2$ needs no training exposure, since
restricting movement reducing contact is knowledge from outside this series. It
behaves as the mechanism predicts and still loses. On the 2020 collapse it helps,
monotonically in $\gamma$ --- 63.45 to 63.28, 63.12 and 62.80, over-prediction
bias falling from $+21.4$ to $+19.9$; on the recovery window straight after it is
destroyed --- 11.54 to 13.45, 15.61 and 18.34, bias flipping from $+0.1$ to
$-6.4$. Pooled it is significantly worse, 30.59, 31.28 and 32.43 against 30.07,
better on 1 of 9 origins at $\gamma \ge 1$ ($p = 0.047$, $p = 0.031$).
Stringency stayed high into 2022 while notifications were already rebounding, so a
multiplier that correctly suppresses transmission in April 2020 goes on
suppressing it for two years too long. A correctly signed, correctly placed,
externally supplied intervention term buys 0.65 RMSE on the window it was built
for and costs 6.81 on the next.

\textbf{What the mechanism is for.} Two roles survive. It constrains: the 39\%
reduction above is obtained by forbidding the network to emit counts that no
compartmental flow could produce. And it supplies quantities a point forecast
cannot --- a transmission rate, an implied $R_{\mathrm{eff}}$, a calibrated
reporting rate and compartment states, one of which is the largest single
contributor to ranking which districts breach an alert threshold
(\S\ref{sec:results}). Neither role asks the mechanism to predict the count, and
that is the distinction this paper draws.

%
% Protocol note: these numbers come from the *corrected* case series (559 weeks,
% `rebuilt`) under nine disjoint origins, not the 459-week array and three
% origins used in the rest of the paper. The first paragraph says so explicitly,
% because a reviewer will otherwise ask. Every figure is reproducible from
% analysis/_build/ and is logged as EXP-032 to EXP-037.
%=============================================================================

\subsection{The mechanism, specified correctly}
\label{sec:mechanism-corrected}

A natural objection to a negative result about mechanistic structure is that the
mechanism was built wrong. We tested that directly. The arm reported here runs on
the date-ordered case series (559 weeks, 25 districts) under nine disjoint
rolling origins spanning 0.50--0.90 of the record, three seeds each; the protocol
is otherwise that of \S\ref{sec:setup}. We state the change of dataset and fold
count plainly because the comparison inside this subsection is internally paired
and should not be read across to Figure~\ref{fig:results}.

Four defects in the compartmental arm were found and corrected, each tested in
isolation: compartments initialised from weekly flows rather than standing stocks,
which caps the attainable prediction at roughly half the observed count;
subtraction of ten years of cumulative cases from the susceptible pool, which
assumes lifelong immunity to a four-serotype disease; a force of infection with
no dependence on current prevalence, so the model contains no epidemic growth;
and a reporting rate fixed at $1/11$ against a published range of $1/2.5$ to
$1/30$. Corrected, the arm improves from test RMSE 50.13 to 29.86, better on
\textbf{9 of 9} origins, $p = 0.0039$ by a sign-flip permutation test clustered
by origin. The improvement replicates on a second construction of the dataset.

\textbf{The compartmental layer earns its place.} The decisive control is the
identical encoder, loss, normalisation and training budget with the simulator
removed from the prediction path, the network predicting counts directly:
\textbf{30.04 with the mechanism against 49.45 without}, better on 8 of 9
origins, $p = 0.0078$ --- a 39\% reduction in error attributable to the structure
alone. It holds at every horizon tested, by 18.2 RMSE at three weeks, 12.1 at six
and 9.0 at twelve, with $p \le 0.012$ in each case. The no-mechanism arm scores
38.23 on validation and 49.45 on test: it fits the periods it has seen and
extrapolates badly, and the compartmental layer removes exactly that freedom.

\textbf{The graph does not.} Freezing the encoder and replacing it with a single
learned transmission rate changes nothing ($p = 0.97$). Replacing border
adjacency with a population-gravity graph does not help either, and the learned
import weight $\alpha$ converges to exactly zero in all 108 runs, for both
graphs: given a free non-negative parameter for its neighbours' prevalence, the
model declines it.

\textbf{And the ceiling holds.} The corrected arm remains level with persistence
rather than ahead of it ($+1.3$ RMSE over nine origins, $p = 0.92$; the best
variant reaches 28.74 against the floor's 28.54). The mechanism for that is
measurable: the corrected forecast correlates with last-value persistence at
$r = 0.97$--$0.99$, differing by 12--15\% of the level, and a learned convex
blend of the two settles at $w = 0.53$--$0.66$ on persistence while changing the
score by 0.04. Trained on squared error, the mechanism rediscovers persistence
and adds a correction worth one to two RMSE, which is the whole margin available.

A fitted $R_{\mathrm{eff}} = \beta S / \gamma$ of $0.40$--$0.44$ explains why.
At a three-week horizon the forecast is dominated by the incubation and reporting
pipeline draining --- people already infected at the origin --- rather than by new
transmission. Two consequences follow, and both are observed. Forecasting further
ahead does not help: persistence stays ahead at six and twelve weeks and the two
degrade at indistinguishable rates, $1.56\times$ against $1.58\times$ from three
weeks to twelve. And an intervention covariate cannot repair the one window
where the record collapses for a non-epidemiological reason. In 2020 national
notifications fall from 3057 to 299 a week under movement restrictions. We gave
the model the OxCGRT stringency index three ways, once with Google workplace
mobility alongside, over 189 further matched evaluations on the same nine origins
and seeds. Learned from the data it cannot act where it is needed: in all 24 runs
whose training data hold no lockdown week both fitted weights stay at exactly
zero, and pooled over nine origins the learned arms are marginally worse than the
cases-only model, 30.05 against 30.06 and 30.07 at Holm $p = 1.0$, by margins
under the pipeline's own run-to-run spread of 0.109 RMSE. Imposing it as a prior
removes that objection: $\beta \mapsto \beta(1 - \sigma/100)^{\gamma}$ with
$\gamma$ fixed in advance at $0.5$, $1$ and $2$ needs no training exposure, since
restricting movement reducing contact is knowledge from outside this series. It
behaves as the mechanism predicts and still loses. On the 2020 collapse it helps,
monotonically in $\gamma$ --- 63.45 to 63.28, 63.12 and 62.80, over-prediction
bias falling from $+21.4$ to $+19.9$; on the recovery window straight after it is
destroyed --- 11.54 to 13.45, 15.61 and 18.34, bias flipping from $+0.1$ to
$-6.4$. Pooled it is significantly worse, 30.59, 31.28 and 32.43 against 30.07,
better on 1 of 9 origins at $\gamma \ge 1$ ($p = 0.047$, $p = 0.031$).
Stringency stayed high into 2022 while notifications were already rebounding, so a
multiplier that correctly suppresses transmission in April 2020 goes on
suppressing it for two years too long. A correctly signed, correctly placed,
externally supplied intervention term buys 0.65 RMSE on the window it was built
for and costs 6.81 on the next.

\textbf{What the mechanism is for.} Two roles survive. It constrains: the 39\%
reduction above is obtained by forbidding the network to emit counts that no
compartmental flow could produce. And it supplies quantities a point forecast
cannot --- a transmission rate, an implied $R_{\mathrm{eff}}$, a calibrated
reporting rate and compartment states, one of which is the largest single
contributor to ranking which districts breach an alert threshold
(\S\ref{sec:results}). Neither role asks the mechanism to predict the count, and
that is the distinction this paper draws.

%
% Protocol note: these numbers come from the *corrected* case series (559 weeks,
% `rebuilt`) under nine disjoint origins, not the 459-week array and three
% origins used in the rest of the paper. The first paragraph says so explicitly,
% because a reviewer will otherwise ask. Every figure is reproducible from
% analysis/_build/ and is logged as EXP-032 to EXP-037.
%=============================================================================

\subsection{The mechanism, specified correctly}
\label{sec:mechanism-corrected}

A natural objection to a negative result about mechanistic structure is that the
mechanism was built wrong. We tested that directly. The arm reported here runs on
the date-ordered case series (559 weeks, 25 districts) under nine disjoint
rolling origins spanning 0.50--0.90 of the record, three seeds each; the protocol
is otherwise that of \S\ref{sec:setup}. We state the change of dataset and fold
count plainly because the comparison inside this subsection is internally paired
and should not be read across to Figure~\ref{fig:results}.

Four defects in the compartmental arm were found and corrected, each tested in
isolation: compartments initialised from weekly flows rather than standing stocks,
which caps the attainable prediction at roughly half the observed count;
subtraction of ten years of cumulative cases from the susceptible pool, which
assumes lifelong immunity to a four-serotype disease; a force of infection with
no dependence on current prevalence, so the model contains no epidemic growth;
and a reporting rate fixed at $1/11$ against a published range of $1/2.5$ to
$1/30$. Corrected, the arm improves from test RMSE 50.13 to 29.86, better on
\textbf{9 of 9} origins, $p = 0.0039$ by a sign-flip permutation test clustered
by origin. The improvement replicates on a second construction of the dataset.

\textbf{The compartmental layer earns its place.} The decisive control is the
identical encoder, loss, normalisation and training budget with the simulator
removed from the prediction path, the network predicting counts directly:
\textbf{30.04 with the mechanism against 49.45 without}, better on 8 of 9
origins, $p = 0.0078$ --- a 39\% reduction in error attributable to the structure
alone. It holds at every horizon tested, by 18.2 RMSE at three weeks, 12.1 at six
and 9.0 at twelve, with $p \le 0.012$ in each case. The no-mechanism arm scores
38.23 on validation and 49.45 on test: it fits the periods it has seen and
extrapolates badly, and the compartmental layer removes exactly that freedom.

\textbf{The graph does not.} Freezing the encoder and replacing it with a single
learned transmission rate changes nothing ($p = 0.97$). Replacing border
adjacency with a population-gravity graph does not help either, and the learned
import weight $\alpha$ converges to exactly zero in all 108 runs, for both
graphs: given a free non-negative parameter for its neighbours' prevalence, the
model declines it.

\textbf{And the ceiling holds.} The corrected arm remains level with persistence
rather than ahead of it ($+1.3$ RMSE over nine origins, $p = 0.92$; the best
variant reaches 28.74 against the floor's 28.54). The mechanism for that is
measurable: the corrected forecast correlates with last-value persistence at
$r = 0.97$--$0.99$, differing by 12--15\% of the level, and a learned convex
blend of the two settles at $w = 0.53$--$0.66$ on persistence while changing the
score by 0.04. Trained on squared error, the mechanism rediscovers persistence
and adds a correction worth one to two RMSE, which is the whole margin available.

A fitted $R_{\mathrm{eff}} = \beta S / \gamma$ of $0.40$--$0.44$ explains why.
At a three-week horizon the forecast is dominated by the incubation and reporting
pipeline draining --- people already infected at the origin --- rather than by new
transmission. Two consequences follow, and both are observed. Forecasting further
ahead does not help: persistence stays ahead at six and twelve weeks and the two
degrade at indistinguishable rates, $1.56\times$ against $1.58\times$ from three
weeks to twelve. And an intervention covariate cannot repair the one window
where the record collapses for a non-epidemiological reason. In 2020 national
notifications fall from 3057 to 299 a week under movement restrictions. We gave
the model the OxCGRT stringency index three ways, once with Google workplace
mobility alongside, over 189 further matched evaluations on the same nine origins
and seeds. Learned from the data it cannot act where it is needed: in all 24 runs
whose training data hold no lockdown week both fitted weights stay at exactly
zero, and pooled over nine origins the learned arms are marginally worse than the
cases-only model, 30.05 against 30.06 and 30.07 at Holm $p = 1.0$, by margins
under the pipeline's own run-to-run spread of 0.109 RMSE. Imposing it as a prior
removes that objection: $\beta \mapsto \beta(1 - \sigma/100)^{\gamma}$ with
$\gamma$ fixed in advance at $0.5$, $1$ and $2$ needs no training exposure, since
restricting movement reducing contact is knowledge from outside this series. It
behaves as the mechanism predicts and still loses. On the 2020 collapse it helps,
monotonically in $\gamma$ --- 63.45 to 63.28, 63.12 and 62.80, over-prediction
bias falling from $+21.4$ to $+19.9$; on the recovery window straight after it is
destroyed --- 11.54 to 13.45, 15.61 and 18.34, bias flipping from $+0.1$ to
$-6.4$. Pooled it is significantly worse, 30.59, 31.28 and 32.43 against 30.07,
better on 1 of 9 origins at $\gamma \ge 1$ ($p = 0.047$, $p = 0.031$).
Stringency stayed high into 2022 while notifications were already rebounding, so a
multiplier that correctly suppresses transmission in April 2020 goes on
suppressing it for two years too long. A correctly signed, correctly placed,
externally supplied intervention term buys 0.65 RMSE on the window it was built
for and costs 6.81 on the next.

\textbf{What the mechanism is for.} Two roles survive. It constrains: the 39\%
reduction above is obtained by forbidding the network to emit counts that no
compartmental flow could produce. And it supplies quantities a point forecast
cannot --- a transmission rate, an implied $R_{\mathrm{eff}}$, a calibrated
reporting rate and compartment states, one of which is the largest single
contributor to ranking which districts breach an alert threshold
(\S\ref{sec:results}). Neither role asks the mechanism to predict the count, and
that is the distinction this paper draws.
